% !TeX root = ../final-report.tex
% TODO(team): Maintain verified prior-work comparisons and citations as the report evolves.
\section{Background and Related Work}
\label{sec:related-work}

\paragraph{Storage and metadata.}
The Google File System (GFS) separates master-managed metadata from
chunkserver data transfers and combines replication, checksums, and
re-replication~\cite{ghemawat2003gfs}. These are established building blocks
for our immutable segment store, rather than new mechanisms. Raft provides
leader election, replicated-log agreement, and state-machine
safety~\cite{ongaro2014raft}. We use it for publication, replica registrations,
and repair-task transitions, not video payloads. Our Controller experiments
test application recovery and quorum-dependent operations; they do not
constitute a new consensus protocol or a proof of storage availability.

\paragraph{Repair timing and interference.}
Total Recall explicitly relates redundancy, transient host absence, repair
timing, and maintenance bandwidth~\cite{bhagwan2004totalrecall}. Its lazy
repair uses additional redundancy to postpone work and avoid unnecessary
copies. Our grace interval studies the same timing concern with a fixed
two-copy starting point, measuring reduced-redundancy exposure instead of
claiming its predicted availability guarantees. In f4, background rebuilders
throttle reconstruction to protect online user
requests~\cite{muralidhar2014f4}. Thus, protecting foreground service from
repair is already established. Our narrower question is how fixed and
source-local load-aware caps change playback stalls and verified replica
restoration under the same total upload budget. Regenerating codes expose a
fundamental storage/repair-bandwidth trade-off~\cite{dimakis2008network};
our prototype uses whole-segment replication and does not implement coding
or claim those bounds.

\paragraph{Streaming and placement.}
CoolStreaming/DONet schedules live-stream segments using availability,
supplier bandwidth, and playback deadlines, and switches suppliers after
departures~\cite{zhang2005coolstreaming}. Our viewers instead read persistent
segments from dedicated storage Peers. This architecture resembles the
distributed-server-network case of Tan and Massouli\'{e}, who distinguish
external requests from requests generated by contributing Peers and analyze
placement for uplink utilization~\cite{tan2011placement}. Zhou et al. model
VoD replication and load balancing, including heterogeneous upload
capacities, to reduce origin-server bandwidth~\cite{zhou2011replication}.
These results establish placement and upload heterogeneity as prior concerns;
their asymptotic or fluid-model conclusions do not establish finite-timeout
playback behavior in our implementation. We therefore hold placement and
round-robin source selection fixed in the main repair experiments.

\paragraph{Playback continuity and scope.}
Zhou et al. analyze the tension between distributing scarce chunks and
fetching urgent chunks for playback~\cite{zhou2011chunks}. BOLA controls
encoded bitrate using buffer occupancy to balance video quality and
rebuffering~\cite{spiteri2016bola}. Our study varies repair service, with
fixed segment payloads and simulated consumption; L observes active
foreground requests rather than playback buffers and is not an adaptive
bitrate algorithm. Together, these papers motivate examining continuity
alongside resource allocation, but do not validate our policy. Our
contribution is a working, verifiable prototype and paired measurements of
playback/repair trade-offs, including conditions where load-aware control
does not consistently improve on fixed caps.
